% !TEX program = pdflatex
% Знако-ранговая карта Шухарта (Shewhart-SR): теория, расчёт контрольных
% границ и разбор предоставленного отчёта.
% Первоисточник: Chakraborti S., Graham M.A. (2019) Nonparametric Statistical
% Process Control, Wiley. Раздел 4.2.1.2, стр. 196--202.
% Сборка: pdflatex (Cyrillic через T2A + babel russian).
\documentclass[11pt,a4paper]{article}

\usepackage[T2A]{fontenc}
\usepackage[utf8]{inputenc}
\usepackage[russian]{babel}
\usepackage[margin=2.4cm]{geometry}
\usepackage{amsmath,amssymb}
\usepackage{pgfplots}
\usepackage{booktabs}
\usepackage{graphicx}
\usepackage{xcolor}
\usepackage{listings}
\usepackage{caption}
\usepackage{hyperref}

\hypersetup{colorlinks=true,linkcolor=blue!50!black,urlcolor=blue!50!black,
            bookmarks=false}

\definecolor{accent}{RGB}{106,76,147}
\definecolor{alert}{RGB}{200,30,30}
\definecolor{softgray}{RGB}{110,110,110}

\lstset{
  basicstyle=\ttfamily\small,
  commentstyle=\itshape\color{softgray},
  keywordstyle=\bfseries,
  breaklines=true,
  frame=single,
  rulecolor=\color{softgray!50},
  columns=fullflexible,
  showstringspaces=false,
  aboveskip=6pt, belowskip=6pt,
}

\pgfplotsset{
  compat=1.18,
  myaxis/.style={
    width=12cm, height=6.6cm,
    tick align=outside,
    axis line style={gray!60},
    tick style={gray!60},
    label style={font=\small},
    title style={font=\small\bfseries},
    every axis plot/.append style={line width=0.9pt},
    legend style={draw=none,font=\small,fill=none},
    grid=none,
  }
}

\newcommand{\SR}{\mathrm{SR}}
\newcommand{\FAR}{\mathrm{FAR}}
\newcommand{\ARL}{\mathrm{ARL}}
\newcommand{\MRL}{\mathrm{MRL}}
\newcommand{\FARz}{\FAR_0}

% ------------------------------------------------------------------ %
%  Данные графиков и таблиц                                           %
% ------------------------------------------------------------------ %
\def\srpmfFive{(-15,0.0312500) (-13,0.0312500) (-11,0.0312500) (-9,0.0625000) (-7,0.0625000) (-5,0.0937500) (-3,0.0937500) (-1,0.0937500) (1,0.0937500) (3,0.0937500) (5,0.0937500) (7,0.0625000) (9,0.0625000) (11,0.0312500) (13,0.0312500) (15,0.0312500)}
\def\srtailfour{(2,0.4375000) (4,0.3125000) (6,0.1875000) (8,0.1250000) (10,0.0625000)}
\def\srtailseven{(2,0.4687500) (4,0.4062500) (6,0.3437500) (8,0.2890625) (10,0.2343750) (12,0.1875000) (14,0.1484375) (16,0.1093750) (18,0.0781250) (20,0.0546875) (22,0.0390625) (24,0.0234375) (26,0.0156250) (28,0.0078125)}
\def\srtailten{(1,0.5000000) (3,0.4609375) (5,0.4228516) (7,0.3847656) (9,0.3476562) (11,0.3125000) (13,0.2783203) (15,0.2460938) (17,0.2158203) (19,0.1875000) (21,0.1611328) (23,0.1376953) (25,0.1162109) (27,0.0966797) (29,0.0800781) (31,0.0654297) (33,0.0527344) (35,0.0419922) (37,0.0322266) (39,0.0244141) (41,0.0185547) (43,0.0136719) (45,0.0097656) (47,0.0068359) (49,0.0048828) (51,0.0029297) (53,0.0019531) (55,0.0009766)}
\def\srtailfourteen{(1,0.5000000) (3,0.4757690) (5,0.4515991) (7,0.4276123) (9,0.4038696) (11,0.3804321) (13,0.3574219) (15,0.3348999) (17,0.3128662) (19,0.2915039) (21,0.2708130) (23,0.2507935) (25,0.2315674) (27,0.2131348) (29,0.1954956) (31,0.1787720) (33,0.1629028) (35,0.1478882) (37,0.1337891) (39,0.1206055) (41,0.1082764) (43,0.0968628) (45,0.0863037) (47,0.0765381) (49,0.0676270) (51,0.0594482) (53,0.0520020) (55,0.0452881) (57,0.0392456) (59,0.0338135) (61,0.0289917) (63,0.0247192) (65,0.0209351) (67,0.0176392) (69,0.0147705) (71,0.0122681) (73,0.0101318) (75,0.0083008) (77,0.0067139) (79,0.0053711) (81,0.0042725) (83,0.0033569) (85,0.0026245) (87,0.0020142) (89,0.0015259) (91,0.0011597) (93,0.0008545) (95,0.0006104) (97,0.0004272) (99,0.0003052) (101,0.0001831) (103,0.0001221) (105,0.0000610)}
\def\srmarksA{(10,0.0625000) (28,0.0078125) (55,0.0009766) (91,0.0011597)}
\def\srmarksSrc{(20,0.0546875) (33,0.0527344) (53,0.0520020)}
\def\srA{(2,3) (3,6) (4,10) (5,15) (6,21) (7,28) (8,36) (9,45) (10,55) (11,64) (12,72) (13,81) (14,91) (15,100) (16,110) (17,121) (18,133) (19,144) (20,156)}
\def\srFAR{(2,0.5000000) (3,0.2500000) (4,0.1250000) (5,0.0625000) (6,0.0312500) (7,0.0156250) (8,0.0078125) (9,0.0039062) (10,0.0019531) (11,0.0019531) (12,0.0024414) (13,0.0024414) (14,0.0023193) (15,0.0026245) (16,0.0026855) (17,0.0025787) (18,0.0023346) (19,0.0023994) (20,0.0023251)}
\def\srSrc{(4,10) (5,15) (6,17) (7,20) (8,24) (9,29) (10,33) (11,38) (12,42) (13,47) (14,53) (15,58)}
\def\srTable{%
2 & 3 & 3 & 0.250000 & 0.500000 & 2.0 & --\\
3 & 6 & 6 & 0.125000 & 0.250000 & 4.0 & --\\
4 & 10 & 10 & 0.062500 & 0.125000 & 8.0 & 10\\
5 & 15 & 15 & 0.031250 & 0.062500 & 16.0 & 15\\
6 & 21 & 21 & 0.015625 & 0.031250 & 32.0 & 17\\
7 & 28 & 28 & 0.007812 & 0.015625 & 64.0 & 20\\
8 & 36 & 36 & 0.003906 & 0.007812 & 128.0 & 24\\
9 & 45 & 45 & 0.001953 & 0.003906 & 256.0 & 29\\
10 & 55 & 55 & 0.000977 & 0.001953 & 512.0 & 33\\
11 & 66 & 64 & 0.000977 & 0.001953 & 512.0 & 38\\
12 & 78 & 72 & 0.001221 & 0.002441 & 409.6 & 42\\
13 & 91 & 81 & 0.001221 & 0.002441 & 409.6 & 47\\
14 & 105 & 91 & 0.001160 & 0.002319 & 431.2 & 53\\
15 & 120 & 100 & 0.001312 & 0.002625 & 381.0 & 58\\
16 & 136 & 110 & 0.001343 & 0.002686 & 372.4 & --\\
17 & 153 & 121 & 0.001289 & 0.002579 & 387.8 & --\\
18 & 171 & 133 & 0.001167 & 0.002335 & 428.3 & --\\
19 & 190 & 144 & 0.001200 & 0.002399 & 416.8 & --\\
20 & 210 & 156 & 0.001163 & 0.002325 & 430.1 & --\\}
\def\srSrcTable{%
4 & 10 & 10 & 0.062500 & 0.1250\\
5 & 15 & 15 & 0.031250 & 0.0625\\
6 & 21 & 17 & 0.046875 & 0.0938\\
7 & 28 & 20 & 0.054688 & 0.1094\\
8 & 36 & 24 & 0.054688 & 0.1094\\
9 & 45 & 29 & 0.048828 & 0.0977\\
10 & 55 & 33 & 0.052734 & 0.1055\\
11 & 66 & 38 & 0.050781 & 0.1016\\
12 & 78 & 42 & 0.054932 & 0.1099\\
13 & 91 & 47 & 0.054932 & 0.1099\\
14 & 105 & 53 & 0.052002 & 0.1040\\
15 & 120 & 58 & 0.053497 & 0.1070\\}
\def\srNfourteenTable{%
105 & 1 & 0.000061 & 0.000061 & 0.000122\\
101 & 1 & 0.000061 & 0.000183 & 0.000366\\
97 & 2 & 0.000122 & 0.000427 & 0.000854\\
95 & 3 & 0.000183 & 0.000610 & 0.001221\\
93 & 4 & 0.000244 & 0.000854 & 0.001709\\
91 & 5 & 0.000305 & 0.001160 & 0.002319\\
89 & 6 & 0.000366 & 0.001526 & 0.003052\\
85 & 10 & 0.000610 & 0.002625 & 0.005249\\
63 & 62 & 0.003784 & 0.024719 & 0.049438\\
61 & 70 & 0.004272 & 0.028992 & 0.057983\\
55 & 99 & 0.006042 & 0.045288 & 0.090576\\
53 & 110 & 0.006714 & 0.052002 & 0.104004\\
51 & 122 & 0.007446 & 0.059448 & 0.118896\\}

\title{\textbf{Знако-ранговая карта Шухарта (Shewhart-SR)}\\[2pt]
\large Теория, расчёт контрольных границ и разбор отчёта с таблицей UCL}
\author{Проект \texttt{nonparametric-charts}}
\date{\today}

\begin{document}
\maketitle

\begin{abstract}
\noindent
Документ описывает знако-ранговую карту Шухарта в случае известной целевой
медианы (Case~K): статистику $\SR_i=2T_n^{+}-n(n+1)/2$, её точное
распределение, правило выбора константы $a$, характеристики длины серии и
примеры. Отдельный раздел разбирает предоставленный отчёт с таблицей
$UCL$ по $n$ и объясняет, почему его значения \emph{не совпадают} с
книжными. Все формулы сверены с Chakraborti \& Graham (2019),
\S\,4.2.1.2, стр.~196--202.
\end{abstract}

\tableofcontents
\bigskip

% ================================================================== %
\section{Назначение и источник}
% ================================================================== %

Знако-ранговая карта (Shewhart-SR) --- непараметрическая карта для
контроля медианы, использующая \emph{и знаки, и ранги} отклонений от
целевого значения. Она мощнее знаковой карты, но требует
дополнительного предположения \emph{симметрии} распределения
относительно $\mu_0$ (Chakraborti \& Graham, 2019, стр.~197).

Первоисточник: Chakraborti, S. \& Graham, M.\,A. (2019).
\textit{Nonparametric Statistical Process Control}. Wiley.
\S\,4.2.1.2 <<Shewhart Control Charts Based on Signed-ranks>>,
стр.~196--202. Расчёт константы $a$ --- стр.~198 (Note~4.2).

% ================================================================== %
\section{Постановка задачи (Case K)}
% ================================================================== %

Подгруппы $X_{i1},\dots,X_{in}$ ($i=1,2,\dots$, $n>1$) --- независимые
наблюдения из неизвестного \emph{непрерывного} распределения,
\emph{симметричного} относительно известного целевого значения $\mu_0$.
Симметрия означает, что среднее и медиана совпадают, поэтому настройка
карты для медианы сопоставима с параметрической картой средних
$\bar X$ (при нормальном распределении).

% ================================================================== %
\section{Статистика}
% ================================================================== %

Пусть $R_{ij}$ --- ранг абсолютного отклонения $|x_{ij}-\mu_0|$ среди
$|x_{i1}-\mu_0|,\dots,|x_{in}-\mu_0|$. Классическая статистика
Уилкоксона (Chakraborti \& Graham, 2019, ур.~(4.4), стр.~197):
\begin{equation}
T_n^{+}=\sum_{j=1}^{n}\Psi(x_{ij}-\mu_0)\,R_{ij},
\end{equation}
а картирующая статистика (ур.~(4.5)):
\begin{equation}
\SR_i=\sum_{j=1}^{n}\phi(x_{ij}-\mu_0)\,R_{ij}
     = \bigl(\text{сумма рангов «плюсов»}\bigr)-\bigl(\text{сумма рангов «минусов»}\bigr).
\label{eq:sr}
\end{equation}
Так как суммы рангов «плюсов» и «минусов» дают в сумме
$S=n(n+1)/2$, получаем линейную связь
\begin{equation}
\SR_i = 2T_n^{+}-\frac{n(n+1)}{2} = 2T_n^{+}-S,
\qquad S=\frac{n(n+1)}{2}.
\label{eq:srrelation}
\end{equation}

\paragraph{Моменты (табл.~4.3, стр.~200).}
В состоянии IC:
\begin{equation}
\mathbb{E}[\SR_i]=0,
\qquad
\operatorname{SD}(\SR_i)=\sqrt{\frac{n(n+1)(2n+1)}{6}} .
\label{eq:srmoments}
\end{equation}

\paragraph{Распределение и чётность.}
Из ур.~\eqref{eq:srrelation}:
$\SR_i=S-2\cdot(\text{сумма рангов «минусов»})$, поэтому $\SR_i$
пробегает значения $S,S-2,\dots,-S$ --- \emph{той же чётности, что и
$S$}. Распределение симметрично относительно нуля. Оно не имеет
замкнутой формы, но вычисляется точно перебором всех $2^n$ знаковых
наборов: каждой из $2^n$ равновероятных комбинаций знаков
$s_j\in\{-1,+1\}$ соответствует значение
$\SR=\sum_{j=1}^{n}s_j\,j$, так что
\begin{equation}
\Pr(\SR=\mathrm{sr})=\frac{\#\{\,s\in\{-1,+1\}^n:\ \textstyle\sum_j s_j j=\mathrm{sr}\,\}}{2^n}.
\label{eq:srpmf}
\end{equation}

% ================================================================== %
\section{Контрольные границы и правило сигнала}
% ================================================================== %

Двусторонние границы и центральная линия
(Chakraborti \& Graham, 2019, ур.~(4.6), стр.~198):
\begin{equation}
\mathrm{UCL}=a,\qquad \mathrm{CL}=0,\qquad \mathrm{LCL}=-a,
\qquad 0<a\le \frac{n(n+1)}{2}=S .
\label{eq:srlimits}
\end{equation}
Сигнал --- если $|\SR_i|\ge a$; пока $-a<\SR_i<a$, процесс в контроле на
уровне $\mu_0$.

% ================================================================== %
\section{Расчёт константы $a$}
% ================================================================== %

\paragraph{Правило (книга, стр.~198).}
$a$ --- \emph{наименьшее} целое $0<a\le S$ такое, что
\begin{equation}
\Pr_{\mu_0\mid IC}(\SR_i\ge a)\;\le\;\frac{\FARz}{2},
\qquad \FARz=0.0027 \;\Rightarrow\; \frac{\FARz}{2}=0.00135 .
\label{eq:srrule}
\end{equation}
Распределение дискретно, поэтому используется \emph{консервативный}
подход (достигнутый FAR не больше номинального). Учитывается чётность:
$a$ должно иметь ту же чётность, что и $S$; иначе значение недостижимо.

\paragraph{FAR и длина серии.}
\begin{equation}
\FAR=\Pr(|\SR_i|\ge a)=2\,\Pr(\SR_i\ge a),
\qquad
\ARL_{IC}=\frac{1}{\FAR},
\qquad
\MRL_{IC}=\left\lceil\frac{\ln 0.5}{\ln(1-\FAR)}\right\rceil .
\label{eq:srarl}
\end{equation}

\paragraph{Fallback при малом $n$.}
Если при данном $n$ вероятность $0.00135$ недостижима, берётся
$a=S$ (самые широкие границы). Так книга поступает в примере~4.4
($n=4$, $a=S=10$).

% ================================================================== %
\section{Примеры по книге}
% ================================================================== %

\paragraph{Пример ($n=5$, стр.~198--199).}
Наименьшая достижимая правосторонняя вероятность равна $0.031>0.00135$;
берётся $T_n^{+}=15$, откуда $a=2\cdot15-15=15$,
$\FAR=2\cdot0.031=0.062$, $\ARL_{IC}\approx16$.

\paragraph{Пример ($n=13$, стр.~199).}
Правосторонний вход $86$: $\Pr(T_n^{+}\ge86)=0.001\le0.00135$, откуда
$a=2\cdot86-91=81$, $\FAR\approx0.002$, $\ARL_{IC}\approx500$.
\emph{Замечание:} книга округляет хвост до $0.001$; точный хвост равен
$0.0012207$, поэтому точный $\FAR=0.002441$ и
$\ARL_{IC}\approx409.6$. Расхождение --- только округление.

\paragraph{Пример 4.4 ($n=4$, стр.~225--228).}
$P(\SR\ge10)=1/16=0.0625>0.00135$, значение недостижимо, поэтому
$a=S=10$, $\FAR=0.125$, $\ARL_{IC}=8$. Книга печатает $\FAR=0.124$
(округление); точное значение $0.125$.

\begin{table}[htbp]\centering
\caption{Контрольная константа $a$ (консервативный подход,
$\FARz=0.0027$) и её сравнение с $UCL$ из предоставленного отчёта
(последний столбец). Точный расчёт.}
\label{tab:srtable}
\begin{tabular}{rrrrrrr}
\toprule
$n$ & $S$ & $a$ & $\Pr(\SR\ge a)$ & $\FAR$ & $\ARL_{IC}$ & $UCL$ отчёта\\
\midrule
\srTable
\bottomrule
\end{tabular}
\end{table}

% ================================================================== %
\section{Графики}
% ================================================================== %

\begin{figure}[htbp]\centering
\begin{tikzpicture}
\begin{axis}[myaxis,
  xlabel={значение $\SR$ при $n=5$},
  ylabel={вероятность},
  xtick={-15,-13,...,15}, ymin=0, ymax=0.12,
  legend pos=north east,
  title={Распределение $\SR$ в состоянии IC ($n=5$)}]
\addplot+[ycomb,mark=*,mark size=2.0pt,color=accent] coordinates {\srpmfFive};
\addlegendentry{$\Pr(\SR=\mathrm{sr})$}
\draw[dashed,alert,line width=1pt] (axis cs:-15,0) -- (axis cs:-15,0.12);
\draw[dashed,alert,line width=1pt] (axis cs:15,0) -- (axis cs:15,0.12);
\end{axis}
\end{tikzpicture}
\caption{Точечная вероятность $\SR$ при $n=5$ (ур.~\eqref{eq:srpmf}),
ср.~рис.~4.3 книги. Здесь $a=S=15$, и хвост $\Pr(\SR\ge15)=1/32$.}
\label{fig:srpmf}
\end{figure}

\begin{figure}[htbp]\centering
\begin{tikzpicture}
\begin{axis}[myaxis,
  xlabel={$a$}, ylabel={$\Pr(\SR\ge a)$},
  ymode=log, log basis y=10,
  xmin=0, xmax=100, ymin=1e-5, ymax=1.3,
  legend pos=south west,
  title={Односторонний хвост $\Pr(\SR\ge a)$}]
\addplot+[mark=o,mark size=1.6pt,color=accent] coordinates {\srtailfour};
\addlegendentry{$n=4$}
\addplot+[mark=square,mark size=1.6pt,color=accent!75!black] coordinates {\srtailseven};
\addlegendentry{$n=7$}
\addplot+[mark=triangle,mark size=1.8pt,color=accent!50!black] coordinates {\srtailten};
\addlegendentry{$n=10$}
\addplot+[mark=diamond,mark size=1.8pt,color=accent!30!black] coordinates {\srtailfourteen};
\addlegendentry{$n=14$}
\addplot[dashed,black,line width=1pt,domain=0:100] {0.00135};
\addlegendentry{книга: $\FARz/2=0.00135$}
\addplot[densely dotted,alert,line width=1pt,domain=0:100] {0.05};
\addlegendentry{отчёт: порог $\approx0.05$}
\addplot[only marks,mark=*,mark size=3.2pt,color=black] coordinates {\srmarksA};
\addlegendentry{книжное $a$}
\addplot[only marks,mark=x,mark size=4pt,color=alert,line width=1.2pt] coordinates {\srmarksSrc};
\addlegendentry{$UCL$ отчёта}
\end{axis}
\end{tikzpicture}
\caption{Односторонний хвост $\Pr(\SR\ge a)$ для разных $n$. Книжное
$a$ (чёрные точки) лежит \emph{ниже} порога $0.00135$; значения отчёта
(красные крестики) --- около $0.05$, то есть примерно в 37 раз выше по
вероятности.}
\label{fig:srtail}
\end{figure}

\begin{figure}[htbp]\centering
\begin{tikzpicture}
\begin{axis}[myaxis,
  xlabel={$n$}, ylabel={граница},
  xmin=1, xmax=21, ymin=0, ymax=160,
  legend pos=north west,
  title={Книжное $a$ против $UCL$ отчёта}]
\addplot+[mark=*,mark size=1.8pt,color=accent] coordinates {\srA};
\addlegendentry{книжное $a(n)$}
\addplot+[mark=x,mark size=2.4pt,color=alert,line width=1.2pt] coordinates {\srSrc};
\addlegendentry{$UCL(n)$ отчёта}
\end{axis}
\end{tikzpicture}
\caption{Книжная константа $a(n)$ (консервативный подход) и границы из
отчёта. Обе монотонно растут, но $UCL$ отчёта систематически ниже,
поскольку соответствует примерно двусторонней вероятности $0.10$, а не
$\FARz=0.0027$.}
\label{fig:srn}
\end{figure}

% ================================================================== %
\section{Реализация (код)}
% ================================================================== %

\begin{lstlisting}
# Exact IC distribution of SR:  sum of +/- j over 2^n sign patterns
def sr_distribution(n):
    counts = {}
    for s in product([-1, 1], repeat=n):
        v = sum(s[j] * (j + 1) for j in range(n))
        counts[v] = counts.get(v, 0) + 1
    return {v: c / 2**n for v, c in counts.items()}

# P(SR >= a) in control
def p_sr_ge_a(n, a):
    return sum(p for v, p in sr_distribution(n).items() if v >= a)

# Conservative charting constant a  (book, p. 198)
def calc_a_conservative(n, far0=0.0027):
    S = n * (n + 1) // 2
    for a in range(1, S + 1):
        if (a - S) % 2 == 0 and p_sr_ge_a(n, a) <= far0 / 2:
            return a, 2 * p_sr_ge_a(n, a)   # (a, FAR)
    return S, 2 * p_sr_ge_a(n, S)           # fallback: widest limits
\end{lstlisting}

% ================================================================== %
\section{Разбор предоставленного отчёта}
% ================================================================== %

\subsection{Что утверждает отчёт}

Ниже воспроизведён фрагмент отчёта. Обозначения: $R^{+}_{tj}$ --- ранг
$|x_{tj}-\theta_0|$; базовая статистика
\begin{equation}
\psi_t=\sum_{j=1}^{n}\operatorname{sign}(x_{tj}-\theta_0)\,R^{+}_{tj},
\qquad
\operatorname{sign}(u)=\begin{cases}+1,&u>0\\0,&u=0\\-1,&u<0,\end{cases}
\end{equation}
связь со статистикой Уилкоксона $W^{+}=\sum_j\Psi(x_{ij}-\theta_0)R_{ij}$:
\begin{equation}
\psi_t = 2W^{+}-\frac{n(n+1)}{2}.
\end{equation}
Отчёт пишет: <<вероятность генерации ложного сигнала (уровень
значимости) $p_0=p(\psi_t\ge UCL)$>>, а границы берутся как
$LCL=-UCL$, и приводит таблицу значений $UCL$ и $p(\psi\ge UCL)$.

\begin{center}\small
\renewcommand{\arraystretch}{1.15}
\begin{tabular}{lcccccccccccc}
\toprule
$n$ & 4 & 5 & 6 & 7 & 8 & 9 & 10 & 11 & 12 & 13 & 14 & 15\\
\midrule
$UCL$ & 10 & 15 & 17 & 20 & 24 & 29 & 33 & 38 & 42 & 47 & 53 & 58\\
$p(\psi\ge UCL)$ & .0625 & .031 & .047 & .055 & .055 & .049 & .053 & .051 & .055 & .055 & .052 & .053\\
\bottomrule
\end{tabular}
\end{center}

\subsection{Что в отчёте верно}

\begin{itemize}
\item \textbf{Статистика совпадает с книжной.}
$\psi_t$ отчёта и $\SR_i$ книги --- это одна и та же величина
($\SR=\psi=2W^{+}-S$), потому что ранг $R^{+}_{tj}$ --- это $R_{ij}$,
а $\operatorname{sign}$ --- это $\phi$. Значит, совпадают и
распределения.
\item \textbf{Связь $2W^{+}-n(n+1)/2$ и симметрия $LCL=-UCL$ --- верны.}
\item \textbf{Числа в столбце $p(\psi\ge UCL)$ точны.} Прямой перебор
всех $2^n$ знаковых наборов воспроизводит их с точностью до округления
(например, $n=4$: $1/16=0.0625$; $n=10$: $\Pr(\psi\ge33)=0.052734$).
То есть арифметика распределения в отчёте правильная.
\end{itemize}

\subsection{Где расхождение: критерий выбора границы}

\emph{Правило книги} (ур.~\eqref{eq:srrule}): $a$ --- наименьшее
достижимое значение с односторонним хвостом
$\Pr(\SR\ge a)\le\FARz/2=0.00135$, то есть под \emph{двусторонний}
риск $\FARz=0.0027$.

\emph{Правило отчёта}: $UCL$ выбирается так, чтобы одностороннее
$p_0=\Pr(\psi\ge UCL)\approx0.05$. Это совершенно другой целевой риск:
$0.05$ примерно в $0.05/0.00135\approx37$ раз больше книжного. Отсюда и
меньшие границы.

\subsection{Одно- и двусторонность: FAR занижен вдвое}

Отчёт называет $p_0=\Pr(\psi\ge UCL)$ <<вероятностью ложного
сигнала>>. Но карта двусторонняя ($LCL=-UCL$): сигнал возникает и при
$\psi\ge UCL$, и при $\psi\le -UCL$. По симметрии истинная вероятность
ложной тревоги равна
\begin{equation}
\FAR=\Pr(|\psi|\ge UCL)=\Pr(\psi\ge UCL)+\Pr(\psi\le -UCL)=2\,p_0 .
\end{equation}
То есть при $p_0\approx0.05$ реальный $\FAR\approx0.10$ --- ложная
тревога примерно каждые $10$ точек. Столбец отчёта даёт \emph{половину}
истинного FAR, и это самостоятельная ошибка трактовки (помимо выбора
$\alpha\ne\FARz/2$).

\subsection{Значения $UCL$ не сопоставимы с книжной константой $a$}

Значения $UCL$ отчёта --- это квантили распределения знаково-ранговой
статистики, подобранные под уровень значимости около $0.05$ (в духе
таблиц критических точек критерия Уилкоксона), пересчитанные по формуле
$\psi=2W^{+}-S$. Они \emph{не} являются <<наименьшим достижимым с
$\Pr\le0.05$>> единообразно (для $n=7$ таким было бы $22$, а в отчёте
$20$; для $n=14$ --- $55$, а в отчёте $53$) и не выводятся из
$\FARz=0.0027$. Поэтому напрямую сравнивать их с книжным $a$ нельзя.

\subsection{Разбор $n=14$}

Для $n=14$: $S=105$, $2^{14}=16384$ наборов, значения $\psi$ нечётные.
Верхний хвост (точный перебор):

\begin{center}\small
\renewcommand{\arraystretch}{1.12}
\begin{tabular}{rrrrr}
\toprule
$\psi$ & число & $\Pr(\psi=\cdot)$ & $\Pr(\psi\ge\cdot)$ & $2\Pr(\psi\ge\cdot)=\FAR$\\
\midrule
\srNfourteenTable
\bottomrule
\end{tabular}
\end{center}

Книжное правило: идём сверху и берём наименьшее $a$ с
$\Pr(\psi\ge a)\le0.00135$. При $a=91$ имеем $0.001160\le0.00135$;
при $a=89$ уже $0.001526>0.00135$. Значит $a=91$,
$\FAR=0.002319$, $\ARL_{IC}\approx431$.

Отчёт даёт $UCL=53$ ($\Pr(\psi\ge53)=0.052$, $2\cdot0.052=0.104$).
Значения $55$ и $63$ на рис.~\eqref{fig:srtail} дополнительно
показывают, что <<наименьшее с $\Pr\le0.05$>> дало бы $55$, а
<<ближайшее к $0.05$ двустороннее>> --- $63$; ни одно из них не равно
$91$.

\subsection{Сравнительная таблица}

\begin{table}[htbp]\centering
\caption{$n=4..15$: книжное $a$ против $UCL$ отчёта.
$\FAR_{отч}=2\Pr(\psi\ge UCL)$ --- истинный двусторонний риск границы
отчёта.}
\label{tab:srsrc}
\begin{tabular}{rrrrr}
\toprule
$n$ & $S$ & $UCL$ отчёта & $\Pr(\psi\ge UCL)$ & $\FAR_{отч}$\\
\midrule
\srSrcTable
\bottomrule
\end{tabular}
\end{table}

\subsection{Вывод}

Отчёт \emph{не} содержит арифметических ошибок в распределении: его
статистика и столбец вероятностей верны. Однако:

\begin{enumerate}
\item \textbf{Неверный целевой риск.} Отчёт проектирует карту под
односторонний уровень $0.05$, тогда как книга --- под
$\FARz/2=0.00135$. Это разные задачи, и значения $UCL$ несопоставимы.
\item \textbf{Занижение FAR вдвое.} $p_0=\Pr(\psi\ge UCL)$ --- это один
хвост; для двусторонней карты $\FAR=2p_0\approx0.10$, а не $0.05$.
\item \textbf{Практическая неприемлемость как книжной карты.} При
$n=14$ граница $53$ даёт $\ARL_{IC}\approx9.6$ --- ложная тревога почти
каждые $10$ точек, что для SPC недопустимо (типичная цель ---
$370$--$500$).
\end{enumerate}

Чтобы получить книжные значения, нужно вернуться к порогу
$\FARz/2=0.00135$ (и к чётности $a\equiv S\pmod2$). Тогда для
$n=4,\dots,15$ получаются
\[
a = 10,\ 15,\ 21,\ 28,\ 36,\ 45,\ 55,\ 64,\ 72,\ 81,\ 91,\ 100
\]
(столбец $a$ табл.~\eqref{tab:srtable}), а не
$10,\ 15,\ 17,\ 20,\ 24,\ 29,\ 33$,
$38,\ 42,\ 47,\ 53,\ 58$.

% ================================================================== %
\section{Итоги}
% ================================================================== %

\begin{itemize}
\item Статистика: $\SR_i=T_n^{+}-T_n^{-}=2T_n^{+}-S$, $S=n(n+1)/2$;
      значения чётности $S$, распределение --- точный перебор $2^n$.
\item Границы: $\mathrm{UCL}=a$, $\mathrm{CL}=0$, $\mathrm{LCL}=-a$;
      сигнал при $|\SR_i|\ge a$.
\item Книжное $a$: наименьшее достижимое с $\Pr(\SR\ge a)\le\FARz/2=0.00135$;
      иначе $a=S$.
\item Отчёт использует порог $\approx0.05$ (в $\sim37$ раз больше) и
      называет односторонний хвост FAR двусторонней карты, занижая риск
      вдвое; поэтому его $UCL$ --- это не книжная константа $a$.
\item $\operatorname{SD}(\SR)=\sqrt{n(n+1)(2n+1)/6}$; $a$ ---
      вероятностный квантиль, а не $3\sigma$.
\end{itemize}

\begin{thebibliography}{9}
\bibitem{chakraborti2019}
S.~Chakraborti, M.\,A.~Graham.
\textit{Nonparametric Statistical Process Control}.
Wiley, 2019. \S\,4.2.1.2, стр.~196--202 (расчёт $a$: стр.~198).
\end{thebibliography}

\end{document}
